\subsubsection{Inviscid\slash Slip Wall}\label{SEC:BC:Wall:Slip}

The inviscid wall boundary condition, sometimes referred to as a slip wall, is
used to model a frictionless solid-wall surface by enforcing the flow velocity
and primitive variable gradients to be tangent to the boundary surface.
%
The details provided here also apply to the symmetry boundary condition
(\ref{SEC:BC:Numerical:Symmetry}) because its formulation happens to be
identical to a slip wall boundary condition.

%%%%%%%%%%%%%%%%%%%%%%%%%%%%%%%%%%%%%%%%%%%%%%%%%%%%%%%%%%%%%%%%%%%%%%%%%%%%%%%%

For this boundary condition, there is no momentum or heat flux through the
boundary so the thermodynamic properties at the wall are copied from the
interior:
%
\begin{alignat*}{1}
  \rho_\subR & = \rho_\subL \\
  T_\subR & = T_\subL
\end{alignat*}
%
To enforce tangency of the velocity at the wall, the wall velocity,
$\vec{u}_\subR$, is computed by taking the interior velocity vector,
$\vec{u}_\subL$, and subtracting out the normal components of the velocity
vector, $(\vec{u}_\subL\cdot\normal)\normal$, resulting in the tangential
velocity at the surface:
%
\begin{equation}
  \vec{u}_\subR = \vec{u}_\subL - \left(\vec{u}_\subL\cdot\normal\right)\normal
  \label{EQ:BC:TangentialVelocity}
\end{equation}
%
For those interested, \eref{EQ:BC:TangentialVelocity} can be
\quoterm{simplified} into a triple vector product as:
%
\begin{equation*}
  \vec{u}_\subR = \normal\times\left(\vec{u}_\subL\times\normal\right)
\end{equation*}
 
%The cross product $\vec{u}\cross\normal$ first creates a vector that is
%perpendicular to the plane created by the velocity and normal vectors.
%%
%Taking the cross product of the normal with the result of the first cross
%product then produces a vector that is perpendicular to the normal and lies in
%the plane created by the velocity and normal vectors, which is just the vector
%tangent to the surface in the direction of the original velocity vector.

%%%%%%%%%%%%%%%%%%%%%%%%%%%%%%%%%%%%%%%%%%%%%%%%%%%%%%%%%%%%%%%%%%%%%%%%%%%%%%%%

For this boundary condition, the wall-normal component of the primitive variable
gradients should be zero at the wall.
%
To accomplish this, we repurpose \eref{EQ:BC:TangentialVelocity} to compute
tangential gradient vectors for each primitive variable at the surface given the
interior gradient vector using
%
\begin{equation*}
  \vec{\varphi}_\subR = \vec{\varphi}_\subL
                   - \left(\vec{\varphi}_\subL\cdot\normal\right)\normal
\end{equation*}
%
where $\vec{\varphi}$ is a stand-in for the gradient of one of the primitive
variables, \thatis \[\vec{\varphi} = \nabla\rho, \nabla u, \nabla v, \nabla w,
\text{ or } \nabla T\]

%###############################################################################

